Electricity demand is rising, driven by electric vehicles, artificial intelligence, and data centers.
To meet this demand, grids draw from various energy sources, each with different carbon factors. 
The combination of these factors is the carbon intensity signal used to calculate the carbon emissions of workloads on the grid. 
The availability of renewable energy sources, weather, and demand all cause the carbon intensity to fluctuate.
Carbon-aware optimizations exploit these fluctuations by shifting workloads to times and places where the grid is cleaner.
These optimizations rely on forecasting models trained to minimize error on large historical datasets, yet the key insight for carbon-aware scheduling is when a region will be greenest, not how green it will be.
We argue that carbon intensity forecasts should target not only accuracy; they must also achieve high \textit{concordance}: the degree to which they preserve the correct ordering of intensity values over time.
We present \textit{LiteCast}, a lightweight time-series forecaster that targets high concordance using only a month of energy, weather, and demand data, which allows it to quickly adapt to grid changes.
Across 48 grids worldwide and real-world workloads, LiteCast delivers 10\% higher carbon savings than state-of-the-art models and achieves 82\% of the maximum carbon savings available for week-long schedules.
